\centerline{\bf \Huge Пробный ЕГЭ I}
\centerline{\bf \Huge }

\begin{flushright}
        \scriptsize{\textbf{76+}}
\end{flushright}

\problem{1} Площадь параллелограмма $ABCD$ равна 20. Точка $E$ — середина стороны $AD$. Найдите площадь трапеции $BCDE$.

\problem{2} Даны векторы $\vec{a}$ и $\vec{b}$. Найдите квадрат длины вектора $\vec{a} - \vec{b}$.

\problem{3} Шар вписан в цилиндр. Площадь полной поверхности цилиндра равна 69. Найдите площадь поверхности шара.

\problem{4} Перед началом волейбольного матча капитаны команд тянут честный жребий, чтобы определить, какая из команд начнёт игру с мячом. Команда «Стартер» по очереди играет с командами «Протор», «Ротор» и «Мотор». Найдите вероятность того, что «Стартер» будет начинать только вторую и последнюю игры.

\problem{5} Игральную кость бросили два раза. Известно, что шесть очков не выпали ни разу. Найдите при этом условии вероятность события «сумма выпавших очков окажется равна 10».

\problem{6} Найдите корень уравнения $2^{4-2 x}=64$.


\problem{7} Найдите значение выражения $\sqrt{48} \cos ^2 \dfrac{19 \pi}{12}-\sqrt{12}$.

\problem{8} На рисунке изображен график $y=f^{\prime}(x)$ - производной функции $f(x)$, определенной на интервале ( $-6 ; 5$ ). В какой точке отрезка $[-5 ;-2]$ функция $f(x)$ принимает наименьшее значение?

\problem{9} В ходе распада радиоактивного изотопа его масса уменьшается по закону $m(t)=m_0 \cdot 2^{-t / T}$, где $m_0$ - начальная масса изотопа, $t$ - время, прошедшее от начального момента, $T$ - период полураспада. В начальный момент времени масса изотопа 128 мг. Период его полураспада составляет 3 мин. Найдите, через сколько минут масса изотопа будет равна 1 мг.

\newpage

\problem{10} Велосипедист выехал с постоянной скоростью из города $A$ в город $B$, расстояние между которыми равно 104 км. На следующий день он отправился обратно со скоростью на 5 км/ч больше прежней. По дороге он сделал остановку на 5 часов. В результате он затратил на обратный путь столько же времени, сколько на путь из $A$ в $B$. Найдите скорость велосипедиста на пути из $A$ в $B$. Ответ дайте в км/ч.


\problem{11} На рисунке изображены графики функций $f(x)=a x^2+b x+c$ и $g(x)=k x+d$, которые пересекаются в точках $A$ и $B$. Найдите абсциссу точки $B$.

\problem{12} Найдите точку минимума функции  $y=(1-5 x) \cos x+5 \sin x+3$  принадлежащую промежутку $\left(0 ; \frac{\pi}{2}\right)$. 

\problem{13} 

a) Решите уравнение $5 \operatorname{tg}^2 x+\dfrac{3}{\cos x}+3=0$.

б) Найдите все корни этого уравнения, принадлежащие отрезку $\left[-\dfrac{5 \pi}{2} ;-\pi\right]$.

\problem{14} В правильной четырёхугольной пирамиде $S A B C D$ сторона основания $A B$ равна 4, а боковое ребро $S A=8$. На рёбрах $C D$ и $S C$ отмечены точки $N$ и $K$ соответственно, причём $D N: N C=S K: K C=1: 3$. Плоскость $\alpha$ содержит прямую $K N$ и параллельна прямой $B C$.

a) Докажите, что плоскость $\alpha$ делит ребро $A B$ в отношении $1: 3$, считая от вершины $A$.

б) Найдите расстояние между прямыми $S A$ и $K N$.

\problem{15} Решите неравенство $1+\dfrac{10}{\log _2 x-5}+\dfrac{16}{\log _2^2 x-\log _2\left(32 x^{10}\right)+30} \geqslant 0$.

\newpage

\problem{16} 15 декабря 2024 года планируется взять кредит в банке на 31 месяц. Условия возврата таковы:

- 1 -го числа каждого месяца долг возрастает на $2 \%$ по сравнению с концом предыдущего месяца;

- co 2-го по 14-е число каждого месяца необходимо выплатить часть долга;

- 15-го числа каждого месяца с 1-го по 30-й (с января 2025 года по июнь 2027 года включительно) долг должен быть на одну и ту же сумму меньше долга на 15 -е число предыдущего месяца;

- 15 июня 2027 года долг составит 100 тысяч рублей;

- 15 июля 2027 года долг должен быть полностью погашен.

Какую сумму планируется взять в кредит, если общая сумма выплат после полного его погашения составит 555 тысяч рублей?

\problem{17} Биссектриса прямого угла прямоугольного треугольника $A B C$ вторично пересекает окружность, описанную около этого треугольника, в точке $L$. Прямая, проходящая через точку $L$ и середину $N$ гипотенузы $A B$, пересекает катет $B C$ в точке $M$.

a) Докажите, $\angle B M L=\angle B A C$

б) Найдите площадь треугольника $A B C$, если $A B=20$ и $C M=3 \sqrt{5}$

\problem{18} Найдите все значения параметра $a$, при каждом из которых система уравнений

$$
\left\{\begin{array}{l}
x+y=2 a \\
|y|=\left|x^2+2 x\right|
\end{array}\right.
$$


имеет два решения.

\problem{19} Дано трёхзначное число $A$, сумма цифр которого равна $S$.

a) Может ли выполняться равенство $A \cdot S=1105$ ?

б) Может ли выполняться равенство $A \cdot S=1106$ ?

в) Какое наименьшее значение может принимать выражение $A \cdot S$, если оно больше 1503 ?


